\documentclass[11pt]{article}
\usepackage[margin=1in]{geometry}
\usepackage{graphicx}
\usepackage{booktabs}
\usepackage{natbib}
\usepackage{hyperref}
\usepackage[T1]{fontenc}
\usepackage[utf8]{inputenc}
\usepackage{amsmath}
\usepackage{microtype}
\sloppy
\setlength{\emergencystretch}{2em}

\title{GaugeFlow: Equivariance, Not Invariance, Is the Right Prior for\\
Continuous Physical Gauges in Self-Supervised Medical Imaging}
\author{Colin Son, MD\\ Seldinger, Inc.\\ \texttt{colinson@gmail.com}}
\date{June 2026}

\begin{document}
\maketitle

\begin{abstract}
A self-supervised objective that works on one medical-imaging task often collapses on the next.
We give a mechanistic account of when it transfers and when it does not, organized around a single
distinction: whether a task's nuisance axis is a \emph{continuous physical gauge} with
\emph{coincident free supervision}, and whether the downstream utility is invariant or covariant
under that gauge. Contrast-front propagation in digital subtraction angiography (DIAS) is such a
gauge: it is a real temporal dynamical process, so ``predict the next state'' is gauge-equivariance
under time-translation with free supervision, and a compact gauge-consistent objective improves
DIAS adjacent-frame prediction while keeping topology-style leakage non-positive. We then ask
whether the same prior transfers to a second continuous physical gauge, the projection angle of
coronary angiography (CardioSYNTAX), where same-study/same-artery retrieval is the utility.
Adversarial invariance via gradient reversal---the standard tool for suppressing a nuisance---is a
no-free-lunch wall here: across a swept adversary strength, no setting holds retrieval non-inferior
while meeting the angle-leakage cap. We test the principled alternative, gauge-equivariance, by
predicting how the embedding should transform under a known $\Delta$angle through an $SO(2)$ head.
As an auxiliary it strictly dominates the adversary, removing its statistically significant harm
($-0.044$, $p\!\approx\!0 \to -0.018$, $p=0.11$) and demonstrably using the gauge (real
$\Delta$angle $0.079$ vs.\ shuffled $0.028$). As the \emph{primary} objective it fails, and the
failure is diagnostic: same-study-across-angles retrieval is itself an invariance metric, and on
this encoder the projection-angle gauge is not a recoverable group action---transport and
shuffled-transport are indistinguishable ($0.0093$ vs.\ $0.0089$) under both common-reference
canonicalization and the rigorous pairwise-transport evaluation. The contribution is a predictive
rule: gauge-equivariant self-supervision helps exactly when the utility is gauge-\emph{covariant}
and the gauge acts as a recoverable transformation of the representation; where the utility is
invariance-shaped, no amount of gauge machinery beats simply not fighting the gauge. All code and
the full per-seed evidence are released.
\end{abstract}

\section{Introduction}

Self-supervised and nuisance-suppression objectives in medical imaging are usually justified by an
intuition---``factor out the scanner,'' ``be invariant to acquisition''---and then validated on one
benchmark. The uncomfortable empirical fact is that the same objective frequently helps on the task
it was tuned for and silently hurts elsewhere. This paper replaces the intuition with a mechanism.
We claim that the deciding property is not the modality or the architecture but a pair of
structural facts about the task: (i) whether the nuisance axis is a \emph{continuous physical
gauge} that comes with coincident, free supervision, and (ii) whether the downstream utility is
\emph{invariant} or \emph{covariant} under that gauge. When both align---a continuous gauge with a
covariant target---a gauge-aware objective is principled and helps. When they do not, the standard
move of suppressing the gauge with an adversary is not merely suboptimal; it is provably caught on
an accuracy-versus-invariance frontier with no free lunch.

We arrive at this through angiography. Contrast-front propagation in DIAS is a genuine temporal
dynamical process: paired frames differ by the gauge (contrast phase) while sharing the same
underlying anatomy-state transition, and ``predict the future observation'' is therefore
gauge-equivariance under time-translation, with the next frame supplying free supervision. A
compact gauge-consistent trainer improves DIAS adjacent-frame prediction over a fixed persistence
comparator while holding topology-style leakage non-positive (Section~\ref{sec:dias}). The
interesting science is what happens on a \emph{second} continuous physical gauge that lacks the
temporal coincidence: the projection angle of an X-ray angiography acquisition.

On that gauge we make three findings. First, adversarial invariance via gradient
reversal~\citep{ganin2016domainadversarial}---the reflexive tool for removing a nuisance---is a
no-free-lunch wall: sweeping its strength traces a monotonic frontier on which no operating point
simultaneously keeps retrieval non-inferior and drives angle-leakage under cap
(Section~\ref{sec:invariance}). Second, the principled alternative suggested by the DIAS
success---gauge-\emph{equivariance}, i.e.\ predicting how the representation should transform under a
known gauge action rather than throwing the gauge away---strictly dominates the adversary when used
as an auxiliary, eliminating its significant harm and provably exploiting the gauge signal
(Section~\ref{sec:equivariance}). Third, promoted to the primary objective, equivariance fails, and
the failure is the point: same-study-across-angles retrieval is structurally an invariance metric,
and the projection-angle gauge is not a recoverable group action on the encoder
(Section~\ref{sec:primary}). We confirm the third finding is protocol-independent with a rigorous
pairwise-transport evaluation, not merely an artifact of how we canonicalize.

The result is a usable rule rather than another benchmark number
(Section~\ref{sec:when}): gauge-equivariant self-supervision helps when, and only when, the task
utility is gauge-covariant and the gauge acts as a recoverable transformation of the representation.
This explains both the DIAS positive and the angiography boundary with one mechanism, and it tells a
practitioner---before training---whether a gauge-aware objective can possibly help on a new task.

\section{The Gauge Frame}\label{sec:frame}

\paragraph{Gauges, invariance, equivariance.}
Let a \emph{gauge} be a controlled, label-preserving change of observation: a transformation $g$ of
the input that leaves the anatomy-state label fixed while changing the appearance. A representation
$z=f(x)$ is gauge-\emph{invariant} if $f(g\cdot x)=f(x)$ and gauge-\emph{equivariant} if there is a
known action $\rho_g$ on representation space with $f(g\cdot x)=\rho_g f(x)$. Invariance discards the
gauge; equivariance keeps it in a structured, predictable form. The two objectives sit on opposite
sides of an information trade: invariance is information-destroying by construction, equivariance is
information-preserving.

\paragraph{Continuous physical gauges with coincident supervision.}
A gauge is a \emph{continuous physical gauge} when it is a real, measurable, continuously
parameterized degree of freedom of the acquisition---contrast phase over time, projection angle in
degrees---rather than a discrete derived label (a style cluster, a vendor bucket). It carries
\emph{coincident supervision} when the gauge parameter is observed for free alongside the data and,
in the strongest case, when moving along the gauge is itself the prediction target. DIAS contrast
propagation is the strong case: the gauge parameter (time) advances and the next frame is the label,
so equivariance under the gauge \emph{is} the task. Projection angle is a continuous physical gauge
with the parameter observed (the positioner angle is recorded) but without temporal coincidence: the
angle is known, yet no free target says ``this is the same anatomy seen $\Delta\theta$ later.''

\paragraph{The covariance condition.}
The deciding question for whether a gauge-aware objective can help is whether the downstream utility
is invariant or covariant under the gauge. If the utility rewards treating gauge-separated views as
the same---e.g.\ retrieving the same study across acquisition angles---then the metric is itself an
invariance objective, and keeping the gauge information works against it. If the utility depends on
the gauge---e.g.\ predicting a gauge-dependent future state---then discarding the gauge destroys
signal the task needs, and equivariance is the principled prior. This single condition, applied
before training, predicts the outcomes in this paper.

\paragraph{The objective.}
For a paired observation $(x_a,x_b)$ of the same cluster differing by gauge value
$(\theta_a,\theta_b)$, with encoder $z=f(x)$, we compare three training signals. \emph{Adversarial
invariance} adds a gradient-reversal head that predicts $\theta_a$ from $z_a$ and is trained to fail,
pushing the gauge out of $z$. \emph{Equivariance (auxiliary)} keeps the base task and adds a head
$h$ that predicts the paired embedding from the gauge displacement,
$\mathcal{L}_{\mathrm{eq}}=\lVert h(z_a,\Delta\theta)-\mathrm{sg}(z_b)\rVert^2$, with
$\Delta\theta=\theta_b-\theta_a$ encoded as $(\sin\Delta\theta,\cos\Delta\theta)$ to respect the
$SO(2)$ structure of a real angle, and a stop-gradient on the target so the contrastive negatives
prevent collapse. \emph{Equivariance (primary)} routes the contrastive positive itself through the
head---transporting $z_a$ into $z_b$'s gauge frame, $h(z_a,\Delta\theta)$, then contrasting against
$z_b$ in both directions---so the only way to satisfy the task is for the gauge to act as a
predictable transformation; at evaluation, views are compared after transport into a common gauge
frame.

\section{Experimental Contract}\label{sec:contract}

Every external number is the output of one pre-registered analyzer: cluster-bootstrap confidence
intervals, a permutation null on the paired per-cluster delta, and a shuffled-gauge negative control
that repartners views at random so the gauge displacement becomes meaningless. A method passes a
gauge only if it holds the task non-inferior to baseline, meets the gauge-leakage cap, and the
shuffled-gauge control reads as null. One portable trainer is used throughout; variants differ only
by configuration flags, so the no-op configuration reproduces the baseline exactly.

The projection-angle gauge uses CardioSYNTAX coronary angiography on real local shards across five
seeds, roughly $200$ study clusters per arm. A precondition for an honest equivariance test is that
the gauge is genuinely continuous; we verified this before training. The recorded primary positioner
angle takes $629$ distinct values over $[-46.4^\circ, 47.5^\circ]$ (standard deviation $19.8^\circ$),
and within-study angle displacements have median $21.7^\circ$ with only $5.4\%$ exactly zero---so the
$11$ coarse buckets used as a categorical gauge label conceal a real continuous signal that an
$SO(2)$ head can exploit. The task is same-study/same-artery retrieval@1 (higher is better); the
leakage probe is the fold-fit $R^2$ with which the angle can be recovered from the representation
(cap $0.035$, the raw-feature reference). The contrast-phase gauge uses the DIAS curated-v1
benchmark against the confirmed persistence comparator.

\section{DIAS: The Gauge That Works}\label{sec:dias}

On the contrast-phase gauge---the one continuous physical gauge in this study whose utility is
covariant, because the task is to predict the gauge-advanced future frame---the gauge-consistent
objective helps. The balanced GaugeFlow variant improves all accepted persistence-baseline metrics:
next-frame test MAE $0.01142$ versus persistence $0.01153$, with a topology-style leakage probe delta
of $-0.074$ (non-positive, the leakage gate) and bounded commutator and state-consistency
diagnostics (Table~\ref{tab:dias}). This is the expected outcome under the gauge frame: time
advancement supplies free supervision, the target is gauge-covariant, and equivariance under
time-translation is exactly what the benchmark rewards.

\begin{table}[h]
\centering
\caption{DIAS contrast-phase gauge (covariant utility). The balanced variant is the first tested
point that clears both the prediction and the topology-style leakage gates.}
\label{tab:dias}
\begin{tabular}{lrrl}
\toprule
Variant & Next-frame MAE & Topology-style leakage $\Delta$ & Gate \\
\midrule
Persistence comparator & 0.011530 & --- & --- \\
GaugeFlow initial & 0.011518 & $+0.021$ & leakage positive, fail \\
GaugeFlow balanced & 0.011418 & $-0.074$ & both gates clear \\
\bottomrule
\end{tabular}
\end{table}

\section{Adversarial Invariance Is a No-Free-Lunch Wall}\label{sec:invariance}

Moving to the projection-angle gauge, the reflexive way to suppress the nuisance is a
gradient-reversal adversary on the embedding. It does not work, and not because it is mistuned. A
single-embedding adversary is actively harmful: retrieval falls from $0.0821$ to $0.0466$, a delta
of $-0.044$ (CI $[-0.063,-0.025]$, permutation $p\approx0$), while angle-leakage stays near baseline
at $R^2=0.31$ against a $0.035$ cap. Porting the full content/style disentanglement of the DIAS
trainer---a dedicated gauge-absorbing path---repairs the destructiveness, cutting angle-leakage
roughly threefold ($0.27\to0.09$) and shrinking the harm to $-0.014$, but it still does not clear the
cap while holding retrieval non-inferior.

Sweeping the adversary strength shows this is a frontier, not a tuning miss (Table~\ref{tab:sweep}).
With the adversary off, the consistency-only objective is non-inferior ($+0.0055$, $p=0.55$) but
leaves leakage at $0.265$; any positive adversary weight drives leakage down at a monotone cost to
retrieval. No swept setting simultaneously holds retrieval non-inferior and meets the leakage cap.
This is the classic invariance information trade made concrete: forcing the gauge out of a
representation that the retrieval task wants to keep coherent necessarily spends task accuracy.

\begin{table}[h]
\centering
\caption{Projection-angle gauge, adversary-strength sweep. The frontier is monotone: invariance
pressure trades retrieval for leakage with no free lunch.}
\label{tab:sweep}
\begin{tabular}{lrr}
\toprule
Adversary weight & Retrieval $\Delta$ vs.\ baseline [CI], $p$ & Angle-leakage $R^2$ \\
\midrule
$0.00$ (consistency only) & $+0.0055$ $[-0.012,+0.024]$, $p=0.55$ & 0.265 \\
$0.05$ (weak) & $-0.044$, harmful & lower \\
$0.25$ (full) & $-0.044$ $[-0.063,-0.025]$, $p\approx0$ & 0.31 \\
\bottomrule
\end{tabular}
\end{table}

\section{Equivariance Dominates Invariance as an Auxiliary}\label{sec:equivariance}

The gauge frame says the principled move is to keep the gauge in a structured form rather than fight
it. We add an $SO(2)$ equivariance head that predicts the paired embedding from the angle
displacement, leaving the retrieval task intact. It strictly dominates the adversary. Retrieval
delta improves from the adversary's significant $-0.044$ ($p\approx0$) to $-0.018$ (CI
$[-0.040,+0.004]$, $p=0.11$): the harm is halved and is no longer statistically distinguishable from
zero. Crucially, the head provably uses the gauge---with the true angle displacement, retrieval is
$0.079$, but with the shuffled-gauge control it drops to $0.028$, a gap of $0.051$ that the adversary
has no analogue for. Leakage stays high ($0.30$) by design, because equivariance keeps the
information rather than destroying it. As an auxiliary, then, equivariance buys back the accuracy the
adversary spends and reads the gauge signal the adversary erases (Table~\ref{tab:frontier}).

\section{Equivariance as Primary Objective: A Diagnostic Failure}\label{sec:primary}

Promoted from auxiliary to primary objective---routing the contrastive positive through the
transport head and retrieving after canonicalization to a common gauge frame---equivariance fails
sharply: retrieval collapses to $0.0089$ against a baseline of $0.095$, a delta of $-0.080$ ($p\approx0$),
with leakage shooting to $0.70$. The failure is informative because the gauge signal vanishes: the
true angle displacement ($0.0089$) is now indistinguishable from the shuffled control ($0.0105$),
the opposite of the auxiliary regime.

Two structural facts explain this. First, same-study-across-angles retrieval treats different angles
of the same study as positives, so the metric \emph{is} an invariance objective; an equivariant
representation that faithfully preserves the angle separates exactly the views the metric wants
together. Second, the projection-angle gauge does not act as a recoverable group action on this
encoder: the transport head cannot map an embedding from one angle frame to another in a way that
recovers study identity, so canonicalization scrambles rather than aligns.

We rule out the obvious objection---that common-reference canonicalization is simply a bad
retrieval protocol---with the rigorous evaluation an equivariant representation deserves: pairwise
transport, moving each query into every candidate's gauge frame before scoring. It gives the same
answer. Retrieval is $0.0093$ for the true gauge versus $0.0089$ for the shuffled control, a null
gap, with leakage $0.73$ (Table~\ref{tab:frontier}). The negative is therefore protocol-independent:
transport and shuffled-transport are indistinguishable however we evaluate, which is direct evidence
that the gauge is not a recoverable transformation of the representation here.

\begin{table}[h]
\centering
\caption{The invariance--equivariance frontier on the projection-angle gauge. Equivariance dominates
adversarial invariance as an auxiliary and uses the gauge; as a primary objective it fails because
the retrieval utility is invariance-shaped and the gauge is not a recoverable group action.
Per-run baselines differ by $\sim\!0.014$ from MPS nondeterminism, so the valid statistic is the
within-run paired delta.}
\label{tab:frontier}
\small
\begin{tabular}{lrll}
\toprule
Objective & Retrieval $\Delta$ vs.\ baseline & Gauge used? (true vs.\ shuffled) & Leakage $R^2$ \\
\midrule
Adversarial invariance (single) & $-0.044$, $p\approx0$ & --- & 0.31 \\
Adversarial invariance (dual path) & $-0.014$ & --- & 0.09 \\
Equivariance, auxiliary & $-0.018$, $p=0.11$ (n.s.) & $0.079$ vs.\ $0.028$ (used) & 0.30 \\
Equivariance, primary (canonical) & $-0.080$, $p\approx0$ & $0.0089$ vs.\ $0.0105$ (none) & 0.70 \\
Equivariance, primary (pairwise) & $-0.050$, $p\approx0$ & $0.0093$ vs.\ $0.0089$ (none) & 0.73 \\
\bottomrule
\end{tabular}
\end{table}

\section{When Do Physical Gauges Help?}\label{sec:when}

The two gauges resolve into one rule. DIAS is a continuous physical gauge with coincident supervision
\emph{and} covariant utility---the task is to predict the gauge-advanced state---so equivariance
under the gauge is the task and the gauge-consistent objective helps. Projection angle is an equally
continuous, equally physical gauge, but its retrieval utility is invariant, and the gauge does not act
as a recoverable transformation of the encoder; there, keeping the gauge fights the metric and
discarding it via an adversary spends accuracy on a no-free-lunch frontier. The best available move
on an invariance-shaped task is the cheapest one---do not fight the gauge at all---which is exactly
why the consistency-only and auxiliary-equivariance settings, both of which leave the task intact,
land at parity while every aggressive gauge objective underperforms.

This yields a pre-training decision rule. Gauge-equivariant self-supervision is the right prior when:
the nuisance is a continuous physical gauge with an observed parameter; the downstream utility is
\emph{covariant} under that gauge; and the gauge acts as a recoverable transformation of the
representation. When the utility is invariance-shaped, no gauge machinery---adversarial or
equivariant---beats not fighting the gauge, and the only objectives that do no harm are the ones that
leave the task alone. The practical payoff is that this is checkable in advance: the continuity of the
gauge is a property of the metadata, the covariance of the utility is a property of the task
definition, and the recoverability of the group action is testable with a transport-versus-shuffle
control before committing to a full training run.

\section{Reproducibility}\label{sec:repro}

The release is self-contained. A single portable trainer implements all four objectives behind
configuration flags---the baseline, the gradient-reversal adversary (single-embedding and dual
content/style path), the auxiliary $SO(2)$ equivariance head, and the primary gauge-aligned
contrastive objective with both common-reference and pairwise-transport evaluation. The
projection-angle dataset adapter reads the CardioSYNTAX shards and exposes the raw positioner angle
as the continuous gauge parameter; the pre-registered analyzer emits every confidence interval,
permutation $p$-value, and shuffled-gauge control in this paper. Each reported number is one
analyzer verdict over five seeds. The external imaging data are licensed and are not redistributed;
the released artifact is the code required to rerun the experiments on the data and the complete
per-seed evidence behind every table.

\bibliographystyle{plainnat}
\bibliography{references}

\end{document}
